%!Mode:: "TeX:UTF-8"

\chapter{绪论}
\echapter{Introduction}  % 英文目录标题

\section{研究背景与意义}
\esection{Research Background and Significance}

学位论文是研究生培养的重要环节，规范的论文格式对于保证学位授予质量具有重要意义。相关背景与排版规范可参见文献\cite{example2024,knuth1984texbook}。

\section{研究内容}

本文的主要研究内容包括：
\begin{enumerate}
  \item 研究内容一；
  \item 研究内容二；
  \item 研究内容三。
\end{enumerate}

\section{公式与图表约定}

本模板的公式按章编号，编号右端对齐，并与公式之间以点线连接（指南三（十））。例如质量守恒的
微分形式可写为

\begin{sduequation}
  \frac{\partial \rho}{\partial t} + \nabla\cdot(\rho\,\mathbf{v}) = 0
  \label{eq:continuity}
\end{sduequation}

式~\eqref{eq:continuity} 中 $\rho$ 为密度，$\mathbf{v}$ 为速度矢量。多行公式仍可使用
\texttt{align} 等环境，编号同样按章编排。

图与表同样按章编号。需要补充说明时使用图注（排在图题下方），用阿拉伯数字按顺序编排
（指南三（八））：

\begin{figure}[htbp]
  \centering
  \includegraphics[width=0.5\textwidth]{figures/SDUWordsColor.pdf}
  \caption{校名标识（示例图）}
  \label{fig:sdu}
  \begin{sdufigurenotes}
    \item 图注位于图题下方，用阿拉伯数字按顺序编号。
    \item 正式论文中请将本图替换为自绘插图。
  \end{sdufigurenotes}
\end{figure}

\section{论文结构}

本文共分为五章，各章内容安排如下：

第一章为绪论，介绍研究背景和意义。

第二章为相关研究，综述国内外研究现状。

第三章为研究方法，详细介绍本文提出的方法。

第四章为实验，验证方法的有效性。

第五章为总结与展望。
